\section*{Use of AI}

This work is the result of human--AI interaction with Claude Opus~5.5 (Anthropic). I chose the problem and directed the project; the AI system found the approach, developed the proofs, and wrote the code, the Lean formalization and the text. Separate AI sessions refereed each proof, checked each citation against its source and recomputed every computational claim with independent code. I take full responsibility for the contents.

\begin{thebibliography}{99}

\bibitem{b1} B. C. Berndt, R. N. Bhat, J. L. Meyer, L. Xie and A. Zaharescu, \emph{An arithmetic sum associated with the classical theta function}, Results Math. \textbf{81} (2026), no.~1, Paper No.~25, doi:10.1007/s00025-025-02584-2; arXiv:2501.03234.

\bibitem{b2} B. C. Berndt, \emph{Analytic Eisenstein series, theta-functions, and series relations in the spirit of Ramanujan}, J. Reine Angew. Math. \textbf{303/304} (1978), 332--365.

\bibitem{b3} B. C. Berndt and L. A. Goldberg, \emph{Analytic properties of arithmetic sums arising in the theory of the classical theta-functions}, SIAM J. Math. Anal. \textbf{15} (1984), 143--150.

\bibitem{b4} S. Bettin, \emph{On the distribution of a cotangent sum}, Int. Math. Res. Not. IMRN \textbf{2015}, no.~21, 11419--11432.

\bibitem{b5} S. Bettin and J. B. Conrey, \emph{Period functions and cotangent sums}, Algebra Number Theory \textbf{7} (2013), 215--242.

\bibitem{b6} S. Bettin and S. Drappeau, \emph{Partial sums of the cotangent function}, J. Th\'eor. Nombres Bordeaux \textbf{32} (2020), 217--230.

\bibitem{b7} S. Bettin and S. Drappeau, \emph{Limit laws for rational continued fractions and value distribution of quantum modular forms}, Proc. Lond. Math. Soc. (3) \textbf{125} (2022), no.~6, 1377--1425.

\bibitem{b8} J. Bourgain and M. Z. Garaev, \emph{Sumsets of reciprocals in prime fields and multilinear Kloosterman sums}, Izv. Math. \textbf{78} (2014), no.~4, 656--707; Russian original: Izv. Ross. Akad. Nauk Ser. Mat. \textbf{78} (2014), no.~4, 19--72; arXiv:1211.4184. Section and page numbers refer to arXiv v1 (18 November 2012), the only arXiv version, unless stated otherwise.

\bibitem{b9} J. Bourgain and M. Z. Garaev, \emph{Kloosterman sums in residue rings}, Acta Arith. \textbf{164} (2014), 43--64.

\bibitem{b10} P. Erd\H{o}s, \emph{On the sum $\sum_{k=1}^{x}d(f(k))$}, J. London Math. Soc. \textbf{27} (1952), 7--15.

\bibitem{b11} Google DeepMind and contributors, \emph{formal-conjectures}, a repository of formalized mathematical conjectures in Lean, \url{https://github.com/google-deepmind/formal-conjectures}, file \texttt{FormalConjectures/Arxiv/2501.03234/ArithmeticSumS.lean} (accessed November 2026).

\bibitem{b12} M. Z. Garaev and I. E. Shparlinski, \emph{On the distribution of modular inverses from short intervals}, Mathematika \textbf{69} (2023), no.~4, 1183--1194.

\bibitem{b13} S. W. Graham and G. Kolesnik, \emph{Van der Corput's Method of Exponential Sums}, London Math. Soc. Lecture Note Ser. \textbf{126}, Cambridge University Press, 1991.

\bibitem{b14} G. H. Hardy, \emph{On certain series of discontinuous functions connected with the modular functions}, Quart. J. Pure Appl. Math. \textbf{36} (1905), 93--123.

\bibitem{b15} K. Henriot, \emph{Nair--Tenenbaum bounds uniform with respect to the discriminant}, Math. Proc. Cambridge Philos. Soc. \textbf{152} (2012), no.~3, 405--424; Erratum, ibid. \textbf{157} (2014), no.~2, 375--377.

\bibitem{b17} M. A. Korolev, \emph{Karatsuba's method for estimating Kloosterman sums}, Sb. Math. \textbf{207} (2016), no.~8, 1142--1158.

\bibitem{b18} A. L\"ageler, \emph{Continued fractions and Hardy sums}, Abh. Math. Semin. Univ. Hambg. \textbf{94} (2024), no.~2, 107--128.

\bibitem{b23} J. B. Rosser and L. Schoenfeld, \emph{Approximate formulas for some functions of prime numbers}, Illinois J. Math. \textbf{6} (1962), 64--94.

\bibitem{b24} P. Shiu, \emph{A Brun--Titchmarsh theorem for multiplicative functions}, J. Reine Angew. Math. \textbf{313} (1980), 161--170.

\bibitem{b25} I. E. Shparlinski, \emph{Modular hyperbolas}, Jpn. J. Math. \textbf{7} (2012), 235--294.

\bibitem{b26} R. Sitaramachandrarao, \emph{Dedekind and Hardy sums}, Acta Arith. \textbf{48} (1987), 325--340.

\bibitem{b27} J. D. Vaaler, \emph{Some extremal functions in Fourier analysis}, Bull. Amer. Math. Soc. (N.S.) \textbf{12} (1985), 183--216.

\bibitem{b34} The mathlib Community, \emph{The Lean mathematical library}, in: Proceedings of the 9th ACM SIGPLAN International Conference on Certified Programs and Proofs (CPP 2020), ACM, 2020, 367--381.

\bibitem{b35} L. de Moura and S. Ullrich, \emph{The Lean 4 theorem prover and programming language}, in: Automated Deduction -- CADE 28, Lecture Notes in Comput. Sci. \textbf{12699}, Springer, 2021, 625--635.

\bibitem{b36} H. D. Kloosterman, \emph{On the representation of numbers in the form $ax^2+by^2+cz^2+dt^2$}, Acta Math. \textbf{49} (1926), 407--464.

\bibitem{b37} R. R. Hall and G. Tenenbaum, \emph{Divisors}, Cambridge Tracts in Mathematics \textbf{90}, Cambridge University Press, 1988.

\bibitem{b38} M. Nair and G. Tenenbaum, \emph{Short sums of certain arithmetic functions}, Acta Math. \textbf{180} (1998), 119--144.

\end{thebibliography}
